\scanpage{76}
Soit $T\subset X$ sous schém\unsure{a|.}, fini étale sur $S$, donc $T\simeq I_S$, $I$ ensemble fini, et soit $U=X\setminus T$ (donc $T$ est défini par une famille $(g_i)_{i\in I}$ de sections disjointes de $X$ sur $S$). \inserted{Si $g$ est le genre relatif, $\nu=\operatorname{card}I$, on suppose $(g,\nu)\neq(0,1)$} Soit \deleted{\unsure{$f_0$|$f$ un|$f$ une}} $i_0\in I$, $f$ une section \deleted{disjointe}\inserted{de $X/S$ distincte} des $g_i$, et telle que \deleted{$f\unsure{\vert|(}K(s)=g_{i_0}$} $f$ et $g_{i_0}$ coïncident en $s$ (pt fermé de $S$). \deleted{\unsure{Considérons}} \inserted{\deleted{\unsure{Donc|Dans}}} \inserted{Si $\eta=S\setminus\{s\}$,} \deleted{\unsure{et|est} le \unsure{groupe|point image|point générique} de} \deleted{par l'homomorphisme\unsure{,}} on a donc un morphisme $\eta\longrightarrow U$, d'où $\Pi_1(\eta)\longrightarrow\Pi_1(U)$. Je dis que cet hom. n'est pas trivial, \deleted{\unsure{donc|(donc}} et $\overline m$, \inserted{\uline{si $\nu\geq2$},} que pour \inserted{$\ell$ donné}
\diagram{\begin{tikzcd}[ampersand replacement=\&,column sep=small]
\& H_1\unsure{(}\eta,\mathbb Z_\ell) \arrow[r] \& H_1(U,\mathbb Z_\ell) \\
\inserted{T_\ell(\unsure{\overline K^*|\widehat K^*})} \arrow[ur,"\unsure{\sim}"] \& \& \inserted{[H_1(U_s,\mathbb Z_\ell\unsure{)}]} \arrow[u,"\sim"']
\end{tikzcd}}
n'est pas trivial (pour, \deleted{\unsure{tout}} $\ell\unsure{=|:=}$ nb premier distinct de la caract\unsure{eris|is}-\newline
tique résiduelle).

Comme la section rationnelle de $\mathcal J^1_{X/S}$ définie par $f$ est régulière, on voit que le composé de l'hom envisagé avec $H_1(\mathcal J^1_{X/S},\mathbb Z_\ell)$ est nul -- i.e. le $H_1(\eta,\mathbb Z_\ell)$ s'envoie dans la partie torique de $H_1(U,\mathbb Z_\ell)$ $\bigl(\xleftarrow{\sim}H_1(U_s,\mathbb Z_\ell)\bigr)$, qui est canoniquement isomorphe à $\deleted{\mathbb Z}\inserted{T}_\ell^I/\deleted{\mathbb Z}\inserted{T}_\ell$.

[NB cette partie est nulle si $\operatorname{card}I=1$, et
\margin{75}
